\section{AI in Meeting Support Systems (RQ2)}

% ---------------------------------------------------------------------------------------------------------------------
\subsection{Overview of existing Meeting AI Tools}
There is a wide variety of AI-based tools that provide support for meetings. Most of these tools focus on
transcribing meeting content. Well-known examples include fireflies.ai \cite{Firefliesai1AI}, Otter.ai
\cite{otteraiOtterMeetingAgent}, and Jamie \cite{jamieaiJamieYourPersonal}.
AI transcription features are also integrated into established meeting platforms such as Google Meet
\cite{workspaceGoogleMeetTelefon}, Zoom \cite{zoomKINotizentoolIhrKIMeetingAssistent},
and Microsoft Teams. \cite{microsoftVideokonferenzenBesprechungenAnrufe}
More advanced features include the automatic extraction of action items. Depending on the tool, these
features are available in real time or only after the meeting has ended.
Microsoft Teams specifically offers a feature called “Facilitator,” \cite{microsoftSetFacilitatorMicrosoft}
which provides additional capabilities beyond the features already mentioned (transcription, identification
of action items): For example, it can edit the agenda if instructed to do so by a user with the appropriate
permissions. It also displays a timeline to track topics and the flow of the meeting, and sends reminders to
manage time. (Once halfway through the meeting and again shortly before it ends.) These interactions take
place via the integrated chat.
Google Meet, for example, also offers a live speech translation feature. \cite{googleLearnSpeechTranslation}

% ---------------------------------------------------------------------------------------------------------------------
\subsection{Research Approaches to AI-Supported Meeting-Systems}
in der litratur finden sich eingie arbieten welche ähnlcihe ansätze verfolgen. so haben z.B. Houtti et al.
einen AI-Agent für inklusivere Meetings desingt. ihre rechere stütze sich auf einem quantitivaen design
interview mit ootenzillen meeting teilneher (nciht wie in dieser arbeit mit facilitating Experten)

In study by Houtti et al.\cite[p. 5]{houttiObserveAskIntervene2025},
in which participants were asked to design virtual co-hosts during semi-structured design sessions,
a widespread tendency emerged to design the co-host to be interactive rather than to let it act entirely
automatically.
Common design patterns involved the co-host providing prompts in response to explicit instructions (e.g.,
when explicitly asked or when a pin was set).
The authors conclude that interactivity could help prevent mistrust.

Acuh ander studien zeigen eine tendenz das wolche systme eher interaktiv gestaltet sein sollten.

\textbf{CoExplorer}

Park et al. \cite{parkCoExplorerTechnologyProbe2024} haben eine adaptiven prototpy entkickler, welcher die
intetinonalität des meetings verbesserl soll. dieser prototyp bestand aus verscheinden Tools, welche die
vorbereitung des Meeting betrafen und die druchfphrung. Teilnehmden hatte so die chance beretis vor dem
meeting mit einem KI-Agenten das Ziel zu sprechen und sich eigen Meinungen abzugene. darauf basierend hat
CoExplorere phasen abgeleite die wahrscheinclih in dem Meeting vorkommen werden. während ees meetings war
Coexplore dann in der Lage die verscheinden Phasen zu erknen und hat über ein adaptives User Interface
relevante ifnroamtionen mit den Teilnehmden live geteilt.

Hierbei fanden fasta lle Teilnehmden an der stuide (24
von 26) das das teilen der Preferneun und meinung im vorraus des mmetings hilfreich ware mich sich tehmatisch
zu alingen und eine priorisierung in der discusions puntk vorzunhmen.

Auch das tracken und visulsieren der Meetingdaten wrude als psotiv aufgenommen.
Einige Teilnehmden außerten sich zu dem system so postiiv, das sie der meinung waren das es keine
menshclichen mediaotr mehr während des meetigns brauche.
Gleichzeitig war aber auch ca. die Hälfte der Teilndhmen der meinung das der Meeting organisator weiterhin in
der verantwortujng bleiben solle und das system nicht proaktiv änderunge vorhenen soll ohne diese mit dem
organizator abgestimmt zu haben, da diese ein besseres verständis über das meeting habe.

% ---------------------------------------------------------------------------------------------------------------------
\subsection{Interaction Paradigms in Human-AI Systems: Reactive, Proactive, and Mixed-Initiative}

% Was das Kapitel abdecken sollte
% 1. Taxonomie der Interaktionsarten (der Kern deiner Frage)

% Reaktiv / responsiv: Agent handelt nur auf explizite Aufforderung (Prompt, Pin, Slash-Command). Das ist der
% Status quo der meisten Tools aus 2.4.1.
% Proaktiv: Agent ergreift selbst die Initiative, ohne explizite Aufforderung.
% Mixed-Initiative: Steuerung wechselt situativ zwischen Mensch und System (Horvitz' Principles of
% Mixed-Initiative User Interfaces ist hier die Schlüsselreferenz).
% Optional: Autonomiegrade / Levels of Automation (Sheridan & Verplank) als Kontinuum statt binärer Entscheidung.
% 2. Abwägung der Trade-offs

% Pro proaktiv: reduziert Prompting-Aufwand, greift Probleme auf, die Teilnehmer selbst nicht ansprechen
% (relevant fürs Facilitating – z.B. stille Teilnehmer, Off-Topic).
% Contra proaktiv: Unterbrechungskosten, Vertrauensverlust, Over-Reliance, Kontrollverlust des Organisators.
% Das verknüpft sich direkt mit deinen Befunden aus 2.4.2 (Houtti et al.: Interaktivität gegen Misstrauen;
% CoExplorer: Hälfte will Organisator in Verantwortung behalten).
% 3. Ableitung / Überleitung

% Schlussfolgerung: Für Meeting-Facilitation ist ein proaktiver, aber kontrollierbarer (mixed-initiative)
% Ansatz angemessen → motiviert direkt 2.4.3 (das Wie des proaktiven Timings).

% ---------------------------------------------------------------------------------------------------------------------
\subsection{Proactive AI-Interactions: From Turn-Taking to instrinsic Motivation with the Inner Thoughts Framework}
To reduce the need for explicit prompts, an AI agent should be able to proactively recognize when it is its
turn in an interaction sequence. In dyadic human-AI interactions, pauses or so-called “stop words” are often
used for this purpose to identify a change in speaker. In multi-party scenarios (e.g., group conversations),
however, this becomes significantly more complex and leads to the “multi-party turn-taking
problem”. \cite{sacksSimplestSystematicsOrganization1974}

Some approaches model the problem as a machine learning task by attempting to predict the next speaker based
on previous speaker changes.
\cite{bayserLearningMultiPartyTurnTaking2019}
Bayer et al. employ traditional algorithms such as Maximum Likelihood Estimation (MLE) and Support Vector
Machines (SVMs), as well as deep learning architectures including Convolutional Neural Networks (CNN) and
Long Short-Term Memory (LSTM) networks. They distinguish two primary paradigms: agent-only
models, which rely solely on speaker history, and agent-and-content models, which additionally incorporate
conversational context. Both paradigms have been trained and evaluated across diverse, topic-oriented and
chit-chat corpora using these varied algorithms.
The empirical findings indicate that while simple agent-only models are sufficient for small, highly focused
dialogues, incorporating conversational content becomes crucial in goal- and topic-oriented interactions
characterized by complex, irregular turn-taking patterns.
\cite[p. 2]{bayserLearningMultiPartyTurnTaking2019}

With these approaches, however, it remains unclear why an agent should interact at a specific point in time.
Due to the design of the ML algorithms, the timing of the interaction is decoupled from its cause or is only
implicitly included in the model, but cannot be explicitly explained. (Example: We know that it is now the
time for the agent to interact now -> But why? Whats the reason its like that he should talk? What should he say?)

Ein naiver ansatz für Turn-taking wäre ein LLM so zu prompnten es zu fragen wer als nächstes dran ist. Liu
et al. \cite{liuProactiveConversationalAgents2025} haben dies überprüft und konnte ihre Hypothese bestätigen,
dass diese Ansatz bei bei turn-allocation ansätzen gut funktioniert. Dies ist nahligen weil in diesem
szenarien in dem chat kontext indizene liegen wer als nächstes dran ist. alle getestetn LLM Modelle hatten
bei turn-acllocation eine besser vorhersagequote als bei selfallocation stellen. die vorhersagen für self
allocation haben sich im rahnen einer zufälligen chance bewegt. Die Authoren schlussfolgern. Das für
self-selection szenarien der konversations kontext allene nicht ausricht und noch intenre factoren dafür
verantwortlich seien.

Die über die leltzten jahre stattgefunden entwicklung bei LLMs ermöglichen neue ansätze um diese Problem zu modellieren.
Neus Paradgima von LLM-basierten Agentischen Systemen.
Was ist das? -> Hier berscheiben was KI-Agenten sind (Anthopic definition finden und verwenden)
LLM basierte AI-Agenten mit einer LLM pipelie welche intrisische gedanken simmuliert, könnte ein viel
versprechender ansatz sein.

Liu et al. verfolgenden diesen Ansatz um einen procativen conversationa Agent mit inner Thoughts zu
modellieren. \cite{liuProactiveConversationalAgents2025}

Y verfolgen diesen ansatz um eine proaktiven Chatbot zu modellieren und
schlussfolgern, dass

\textbf{Classic Worksflow}

% ---------------------------------------------------------------------------------------------------------------------
% \subsection{Theoretical frameworks: Human-AI Interaction (HAI), Explainable AI (XAI) }
